\documentclass[12pt]{article}
\usepackage[utf8]{inputenc}
\usepackage[T1]{fontenc}
\usepackage{hyperref}
\title{Coatecas Altas Zapotec Language Resources for K-12 Curriculum Development}
\author{Compiled by Claude Sonnet 5.5 (claude-sonnet-5-5) for Andr\'as Kornai}
\date{}
\begin{document}
\maketitle

\section{Introduction}
Coatecas Altas Zapotec (ISO 639-3: zca; Glottocode coat1244) is an Otomanguean language of the Zapotec macrolanguage, spoken in and around Coatecas Altas in the Ejutla District of southern Oaxaca, Mexico. A speaker count of about 4,900 (2005 census) is cited. Some speakers also live in Chiapas, in San Bernardino (Oaxaca) and near San Quint\'in, Baja California. Wikitongues rates it EGIDS 6a (vigorous). The web search found very little material beyond catalogue entries and religious materials. No school textbook, primer, dictionary or grammar turned up.

\section{Primary Educational Resources}
No primers, textbooks or school programmes for this language were found in search results. Any curriculum would have to be built from scratch, probably starting from the scripture texts and recordings listed below.

\section{Linguistic and Reference Materials}
\begin{itemize}
\item Wikipedia entry (English), with basic classification and speaker data: \url{https://en.wikipedia.org/wiki/Coatecas_Altas_Zapotec}
\item Wikitongues language page, with the EGIDS vitality rating: \url{https://wikitongues.org/languages/zca/}
\item Glottolog bibliographic entry (Joseph P. Benton Scott 2016) that appeared in search results; its content was not verified: \url{https://glottolog.org/resource/reference/id/552060}
\end{itemize}

\section{Religious and Community Resources}
\begin{itemize}
\item ScriptureEarth page for zca, listing a Bible in the language with read, listen and view options, a New Testament PDF, and mobile apps: \url{https://www.scriptureearth.org/00eng.php?iso=zca}
\end{itemize}

\section{Digital and Audio Resources}
\begin{itemize}
\item Global Recordings Network page for zca, with audio-visual Bible lessons (Good News series): \url{https://globalrecordings.net/en/language/zca}
\end{itemize}

\section{Curriculum Development Context}
Little published material exists for zca in the sources found. Scripture translations and audio recordings are the only substantial texts, and no dictionary, grammar or orthography guide was located. Further work would need to check OLAC and similar catalogues, and ideally consult speaker communities in Coatecas Altas. Related Zapotec varieties (for example Ozolotepec and Miahuatl\'an) may offer comparative material, but they are different languages and are not covered here.
\end{document}